\documentclass[11pt]{report}
\usepackage{mystyle}
\usepackage{epigraph}

\begin{document}

%% TITLE PAGE %%
\thispagestyle{empty}
\begin{center}

%    \newgeometry{top=1cm,bottom=1cm,outer=1cm,inner=1cm}
    \rule{10cm}{0.5mm}

    \huge
    {\sc
        Ab Initio Modelling of \\
        Photoinduced Electron Dynamics \\
        in Nanostructures 
    }

    \vspace{-0.5cm}
    \rule{10cm}{0.5mm}

    \vspace{1.5cm}
    \Large 
    A Thesis \\[0.5cm]
    Presented Upon Application for \\
    Admission to the Degree of \\[0.5cm]

    \LARGE
    {\sc Doctor of Philosophy} \\[0.5cm]

    \Large 
    in the\\
    Faculty of Engineering and Physical Sciences\\[0.5cm]
    by \\[0.2cm]

    \LARGE
    \textbf{Ryan J. McMillan} \\
    MSci (Hons) 2013 \\[0.7cm]

    \includegraphics[height=4cm]{img/QUB.eps}

    \vspace{0.5cm}
    \Large
        School of Mathematics and Physics\\
        Queen's University Belfast\\
        Northern Ireland\\[0.5cm]
    
    \LARGE
    \monthyeardate\today

    \normalsize
\end{center}

%%

\pagenumbering{roman}

%% TABLE OF CONTENTS %%
%\setstretch{1.5}
\clearpage
\tableofcontents

%% LIST OF FIGURES %%
\clearpage
\listoffigures
%%

\pagenumbering{arabic}
\setlength{\topmargin}{0.4cm}
\setlength{\headsep}{0.6cm}
%\setstretch{1.5}

%% INTRO %%
\clearpage
\chapter{Background and Introduction}

\section{Motivation - Solar Technology}
A brief discussion about renewable energy and solar technology in particular.
We'll mention the various types of latest generation solar cells with a focus
on 2D materials.
\section{2D Materials in Photonics}
Properties of 2D materials that make them useful in photonics with
experimental and theoretical results.
\section{Increasing Light Absorption in 2D Materials}
Different methods of increasing light absorption with attention towards
plasmonic field enhancement.
%\section{Optical Spectra for \texorpdfstring{\ce{MoS2}}{MoS2} Composites}\label{sec:mos2}

%\input{tex/intro/big_pic.tex}
%\input{tex/intro/small_pic.tex}
%\input{tex/intro/nano_pic.tex}
%%

%% THEORY %%
\clearpage
\chapter{Theory}

\section{Background Theory}
\subsection{Electronic Structure Calculations}
Brief discussion of various methods for calculating electronic structure,
focussing on Kohn-Sham DFT.
\subsection{Excited State Calculations}
Brief discussion of TDDFT, linear response and the (time-dependent)
Bethe-Salpeter Equation.
\subsection{Obtaining Optical Properties}
A bit about dielectric functions and absorption spectra and how they are
obtained from excited state calculations.
\section{Projected Equations of Motion Method for Composite Materials}
This will include a little about obtaining the polarization from classical
electrodynamics. The dipole-dipole model will be presented and the PEOM method
will be derived.
%\section{Optical Spectra for \texorpdfstring{\ce{MoS2}}{MoS2} Composites}\label{sec:mos2}

%\input{tex/theory/elec_struc/ipm.tex}
%%

%% RESULTS %%
\clearpage
\chapter{Results}

Here we use the SQD-MNP system to validate the method as well as providing new
results for femtosecond pulse interaction. The method is further verified by
performing calculations on various monolayer systems using TDDFT and BSE
theory. We'll then go on to model a monolayer of \mos with gold nanoparticles.

\section{Semiconducting Quantum Dot-Metal Nanoparticle Hybrid}
\subsection{Describing the System}
\subsection{Energy Absorption Rates}
\subsection{Population Inversion}
\section{Optical Spectra for \mos Composites}
\subsection{\mos Monolayer}
\subsection{\mos / Graphene Monolayers}
\subsection{\mos Bilayer}
\subsection{\mos Monolayer on Gold Surface}
\subsection{\mos Monolayer with Gold Nanoparticles}
%%

%% CONCLUSION %%
\clearpage
\chapter{Conclusions}
We will have shown that the PEOM method provides a way for modelling the
plasmonic field enhancement from gold nanoparticles on a \mos monolayer. Such
theoretical calculations have not been performed before.

\setstretch{1.0}
%% BIBLIOGRAPHY %%
\clearpage
\bibliography{$HOME/Dropbox/PhD/latex-files/Master,$HOME/Dropbox/PhD/latex-files/websites}
%%

\end{document}
